%=============================================================
% Response to Reviewers — Hyperparameter Selection Defense
% Paper: Addressing the Long-Tail in Cold-Start DDI Prediction
%        via Graph Attention and Fingerprint Fusion
% Author: Buddharak Sodtana
%=============================================================
\documentclass[11pt,a4paper]{article}

\usepackage[utf8]{inputenc}
\usepackage[T1]{fontenc}
\usepackage[top=1in, bottom=1in, left=1.1in, right=1.1in]{geometry}
\usepackage{amsmath,amssymb}
\usepackage{booktabs}
\usepackage{xcolor}
\usepackage{mdframed}
\usepackage{enumitem}
\usepackage{hyperref}
\usepackage{parskip}
\usepackage{titlesec}
\usepackage{setspace}

\hypersetup{colorlinks=true, linkcolor=blue, urlcolor=blue}

% ---- colours ----
\definecolor{reviewerbg}{RGB}{245,245,250}
\definecolor{reviewerbar}{RGB}{90,100,180}
\definecolor{authorbg}{RGB}{250,250,245}
\definecolor{authorbar}{RGB}{60,150,90}
\definecolor{warningbg}{RGB}{255,248,230}
\definecolor{warningbar}{RGB}{210,150,30}

% ---- reviewer comment box ----
\newmdenv[
  backgroundcolor=reviewerbg,
  linecolor=reviewerbar,
  linewidth=3pt,
  leftline=true, rightline=false, topline=false, bottomline=false,
  innerleftmargin=10pt, innerrightmargin=8pt,
  innertopmargin=6pt, innerbottommargin=6pt,
  skipabove=8pt, skipbelow=4pt
]{reviewerbox}

% ---- author response box ----
\newmdenv[
  backgroundcolor=authorbg,
  linecolor=authorbar,
  linewidth=3pt,
  leftline=true, rightline=false, topline=false, bottomline=false,
  innerleftmargin=10pt, innerrightmargin=8pt,
  innertopmargin=6pt, innerbottommargin=6pt,
  skipabove=4pt, skipbelow=8pt
]{authorbox}

% ---- risk/caution box ----
\newmdenv[
  backgroundcolor=warningbg,
  linecolor=warningbar,
  linewidth=3pt,
  leftline=true, rightline=false, topline=false, bottomline=false,
  innerleftmargin=10pt, innerrightmargin=8pt,
  innertopmargin=6pt, innerbottommargin=6pt,
  skipabove=4pt, skipbelow=8pt
]{cautionbox}

\titleformat{\section}{\large\bfseries}{}{0em}{}[\titlerule]
\titlespacing*{\section}{0pt}{14pt}{6pt}
\titleformat{\subsection}{\normalsize\bfseries\itshape}{}{0em}{}
\titlespacing*{\subsection}{0pt}{10pt}{3pt}

\newcommand{\reviewer}[1]{\textbf{\textcolor{reviewerbar}{Reviewer:}} #1}
\newcommand{\authorlabel}{\textbf{\textcolor{authorbar}{Response:}}}

\setstretch{1.10}
\setlist{nosep, leftmargin=*}

%=============================================================
\begin{document}
%=============================================================

\begin{center}
  {\LARGE\bfseries Response to Reviewers}\\[0.5em]
  {\large Addressing the Long-Tail in Cold-Start Drug--Drug Interaction\\
  Prediction via Graph Attention and Fingerprint Fusion}\\[0.5em]
  {\normalsize Buddharak Sodtana \quad|\quad
  Advisor: Assoc.\ Prof.\ Prompong Sugunnasil}\\[0.3em]
  {\small\textcolor{gray}{Manuscript revised March 2026}}
\end{center}

\medskip
\noindent\rule{\textwidth}{0.4pt}
\medskip

\noindent\textit{%
We thank the reviewers for their careful reading and constructive comments.
Below we address each point in turn.
Reviewer comments appear in blue-bordered boxes; our responses in green-bordered boxes.
All section and table references are to the revised manuscript.}

%=============================================================
\section{Reviewer Comment C: Non-Standard Notation in Equation~3}
%=============================================================

\begin{reviewerbox}
\reviewer{In Equation~3, you define the pair feature as
$[\mathbf{h}_{\min} \| \mathbf{h}_{\max} \| (\mathbf{h}_{\max} - \mathbf{h}_{\min}) \| \mathbf{h}_a \odot \mathbf{h}_b]$.
The term $\mathbf{h}_{\max} - \mathbf{h}_{\min}$ is simply the absolute difference
$|\mathbf{h}_a - \mathbf{h}_b|$.
Write it as the absolute difference.
The current notation is unnecessarily convoluted and implies a directional
subtraction that does not exist.}
\end{reviewerbox}

\begin{authorbox}
\authorlabel\ We agree.
Since $\mathbf{h}_{\min}$ and $\mathbf{h}_{\max}$ are defined element-wise,
$\mathbf{h}_{\max} - \mathbf{h}_{\min}$ is indeed equal to $|\mathbf{h}_a - \mathbf{h}_b|$
component by component.
Using the difference-of-sorted-extrema form was needlessly indirect and could
suggest a directed subtraction that has no meaning for an unordered drug pair.

In the revised manuscript, Equation~3 now reads:
\begin{equation}
  \mathbf{z}_{ab} =
  [\mathbf{h}_{\min} \| \mathbf{h}_{\max} \|
   |\mathbf{h}_a - \mathbf{h}_b| \|
   \mathbf{h}_a \odot \mathbf{h}_b],
\end{equation}
where $|\cdot|$ denotes the element-wise absolute difference.
A one-line bridging statement is also added in the text:
``Note that $\mathbf{h}_{\max} - \mathbf{h}_{\min} = |\mathbf{h}_a - \mathbf{h}_b|$
element-wise,'' so readers can confirm the equivalence without consulting the code.
\end{authorbox}

%=============================================================
\section{Reviewer Comment: Hyperparameter Selection (\S3.7, Table~2)}
%=============================================================

\begin{reviewerbox}
\reviewer{The hyperparameter values in Table~2 (e.g., $\tau = 0.5$,
learning rate $10^{-3}$, DRW start epoch 15) appear without justification.
How were these chosen?
Were they tuned on the test set?}
\end{reviewerbox}

\begin{authorbox}
\authorlabel\
We agree that the hyperparameters in Table~2 should not appear as unexplained
``magic numbers.''
In the revised manuscript we now state the selection procedure explicitly.

The reported values were \emph{not} chosen from the full five-fold benchmark results
and were \emph{not} tuned on the held-out test folds.
Instead, we used a bounded development-fold search implemented in
\texttt{scripts/tune\_baseline\_v4.py}.
This search fixed the model family (\texttt{GAT + ECFP}) and screened a predefined
grid on a single cold-drug development fold (fold~0), using validation performance
only for selection.

The screened search space was:
\begin{itemize}
  \item Logit-adjustment temperature
        $\tau \in \{0.4,\;0.5,\;0.75\}$
  \item Learning rate
        $\in \{3\!\times\!10^{-4},\;5\!\times\!10^{-4},\;10^{-3},\;2\!\times\!10^{-3}\}$
  \item DRW ratio $\in \{0.6,\;0.7,\;0.8\}$, with learning-rate drop $= 0.2\times$
\end{itemize}

This corresponds to 36 planned configurations; 31 completed runs are recorded in
the saved sweep export
(\texttt{outputs/tuning\_baseline\_v4\_single\_fold/sweep\_results\_31.json}).
Configurations were ranked by the pre-declared priority order:
\[
  \text{tail macro PR-AUC}
  \succ \text{macro PR-AUC}
  \succ \text{macro-F1}
  \succ \text{lower validation objective loss.}
\]

Under this procedure, the top-ranked configuration was $\tau = 0.5$,
learning rate $10^{-3}$, and DRW ratio $0.7$ (DRW start epoch~15).
The runner-up configuration with DRW ratio $0.8$ tied on the primary validation
criteria and was also shortlisted.

Architectural values such as 3 GAT layers, hidden dimension 64, output dimension 128,
4 attention heads, and dropout~0.2 were held fixed during this search.
These were treated as controlled baseline settings rather than exhaustively
optimized architecture parameters.
Our intent was to build a reproducible cold-drug benchmark under a fixed compute
budget and a fixed optimization protocol, not to claim globally optimal tuning
across all possible architectures.

We therefore describe the study as using a \emph{structured but bounded}
hyperparameter search.
This is stronger than ad hoc trial-and-error, but it is not a full nested
hyperparameter optimization study.
We have clarified this limitation in the manuscript (revised \S3.7).
\end{authorbox}

%---------------------------------------------------------------
\subsection{Follow-up: Why not tune on full data?}
%---------------------------------------------------------------

\begin{reviewerbox}
\reviewer{Why not perform hyperparameter tuning on the full dataset to obtain
stronger baselines?}
\end{reviewerbox}

\begin{authorbox}
\authorlabel\
We did not tune on the full benchmark because that would mix model selection
with final evaluation.
The correct protocol is to tune on a development split and keep the outer test
folds untouched.
Our choice of a single development fold plus a 100{,}000-row training cap was a
compute-control decision, not an attempt to weaken baselines.
Using the benchmark test folds for tuning would constitute evaluation leakage and
inflate reported performance---precisely the methodological flaw we aim to
correct relative to prior warm-start DDI studies.
\end{authorbox}

%---------------------------------------------------------------
\subsection{Follow-up: Why were 3 layers and 128-dim not searched?}
%---------------------------------------------------------------

\begin{reviewerbox}
\reviewer{The paper does not justify the choice of 3 GAT layers and 128-dimensional
embeddings.
Were these also swept?}
\end{reviewerbox}

\begin{authorbox}
\authorlabel\
Those architecture parameters were fixed \emph{a priori} for the controlled
benchmark and were not part of the swept search space.
We will revise the text to state that clearly rather than presenting them as if
they were exhaustively optimized.

The rationale for fixing architecture dimensions is reproducibility and compute
control: running a full joint search over architecture $\times$ optimization
hyperparameters would require nested cross-validation, which is infeasible at the
scale of the 86-class cold-drug benchmark within a single workstation budget.
This is an acknowledged limitation; we recommend a full nested HPO study as
priority future work (revised \S5, Future Work).
\end{authorbox}

%---------------------------------------------------------------
\subsection{Revised methodology paragraph (inserted in \S3.7)}
%---------------------------------------------------------------

The following paragraph has been added to \S3.7 (Experimental Configuration)
of the revised manuscript:

\begin{mdframed}[backgroundcolor=gray!8, linewidth=0pt,
                 innerleftmargin=12pt, innerrightmargin=12pt,
                 innertopmargin=8pt, innerbottommargin=8pt]
\noindent\textit{Hyperparameter selection.}
The final configuration was chosen by a bounded development-fold search rather
than by tuning on the full five-fold benchmark.
Specifically, we used a predefined single-fold screening procedure on cold fold~0
with a training cap of 100{,}000 rows for compute efficiency, using
validation-based model selection only.
The search space was:
logit-adjustment temperature $\tau \in \{0.4,\,0.5,\,0.75\}$,
learning rate $\in \{3\!\times\!10^{-4},\;5\!\times\!10^{-4},\;10^{-3},\;2\!\times\!10^{-3}\}$,
and deferred re-weighting ratio $\in \{0.6,\,0.7,\,0.8\}$ with
learning-rate drop fixed at $0.2\times$.
Configurations were ranked by validation tail macro PR-AUC, then macro PR-AUC,
then macro-F1, and finally lower validation objective loss.
This screening selected $\tau=0.5$, learning rate $10^{-3}$, and DRW ratio~$0.7$
(start epoch~15).
Architectural settings such as 3 GAT layers, hidden dimension~64, output
dimension~128, 4 attention heads, and dropout~0.2 were held fixed during this
search to maintain a controlled benchmark under a fixed compute budget.
No outer-fold test set was used for hyperparameter tuning.
\end{mdframed}

%=============================================================
\section{Concrete Evidence for Hyperparameter Procedure}
%=============================================================

The following artifacts are available in the repository to verify the procedure:

\begin{table}[h]
\centering
\small
\begin{tabular}{ll}
\toprule
\textbf{Artifact} & \textbf{Path} \\
\midrule
Search implementation  & \texttt{scripts/tune\_baseline\_v4.py} \\
Raw sweep table        & \texttt{outputs/tuning\_baseline\_v4\_single\_fold/screen\_results\_raw.csv} \\
Ranked export (31 runs)& \texttt{outputs/tuning\_baseline\_v4\_single\_fold/sweep\_results\_31.json} \\
Shortlisted top-2 configs & \texttt{scripts/benchmark\_v6\_top2.sh} \\
\bottomrule
\end{tabular}
\caption{Repository artifacts documenting the hyperparameter search procedure.}
\end{table}

\medskip
\noindent Key numbers if challenged:

\begin{table}[h]
\centering
\small
\begin{tabular}{ll}
\toprule
\textbf{Item} & \textbf{Value} \\
\midrule
Planned grid size       & 36 configurations \\
Completed saved runs    & 31 \\
$\tau$ values swept     & $\{0.4,\;0.5,\;0.75\}$ \\
Learning rates swept    & $\{3\!\times\!10^{-4},\;5\!\times\!10^{-4},\;10^{-3},\;2\!\times\!10^{-3}\}$ \\
DRW ratios swept        & $\{0.6,\;0.7,\;0.8\}$ \\
Ranking rule            & tail PR-AUC $\succ$ macro PR-AUC $\succ$ macro-F1 $\succ$ lower loss \\
Top-ranked config       & $\tau=0.5$, lr $= 10^{-3}$, DRW ratio $= 0.7$, start epoch~15 \\
Runner-up (tied on dev) & $\tau=0.5$, lr $= 10^{-3}$, DRW ratio $= 0.8$, start epoch~17 \\
\bottomrule
\end{tabular}
\caption{Exact sweep parameters and results for reviewer verification.}
\end{table}

%=============================================================
\section{Risks and Honest Limitations}
%=============================================================

\begin{cautionbox}
\textbf{Caution --- What not to claim in the response:}
\begin{itemize}
  \item \textbf{Do not} claim ``rigorous full hyperparameter optimization.''
        The code does not support this: only loss/schedule parameters were swept,
        architecture was fixed, training was capped at 100k rows, and only one
        development fold was used.
  \item \textbf{Do not} say ``no tuning was done.''
        That is also false: a real grid search was conducted and documented.
\end{itemize}

\medskip
\noindent\textbf{The honest position is:}
\begin{itemize}
  \item Structured search for loss/schedule hyperparameters
  \item Fixed architecture (controlled baseline, not optimized)
  \item Single development fold (not nested cross-validation)
  \item Compute-capped training subset (100k rows)
\end{itemize}

\medskip
\noindent\textbf{If the reviewer insists on the strongest possible standard},
the recommended next experiment is:
\begin{enumerate}
  \item Full nested HPO on inner validation only (architecture + optimization search)
  \item Freeze the winner
  \item Rerun all five outer test folds on the frozen configuration
\end{enumerate}
\end{cautionbox}

\end{document}
